%% Formule de l'operateur : PressAltiOACI
%% Dans ce fichier, les �ventuels accents sont cod�s
%% en ISO-8859-15 (aka latin9)
%% La ligne suivante donne trois exemples 
%% e accent aigu : �, e accent grave : �, e-dans-l'o : �


Altitude g�opotentielle $H= (\frac{r*h}{r+h}) $\\
Rayon de la terre r=6356766,0 m\\
Acc�leration normalis�e de la pesanteur  $g_0$ = 9,80665 $m/s^2$ \\
Constante massique des gaz parfaits sp�cifique de l'air  R  = 287,05287 $J/(K*kg)$\\

Si H inf�rieure � -5000 m  \\
 hors de la norme OACI, $P=misval$\\

Si H comprise entre -5000 et 0 m  \\
$P = P_b[1+ \frac{\beta }{T_b}(H-H_b)]^{(\frac{-g_0}{\beta R})}$\\
avec $P_b$ = 1776,87 hPa\\
et $\beta$ = -0,0065 K/m\\
et  $T_b$ = 320,65 K\\
et  $H_b$  = -5000,0 m\\

Si H comprise entre 0 et 11000 m  \\
$P = P_b[1+ \frac{\beta }{T_b}(H-H_b)]^{(\frac{-g_0}{\beta R})}$ \\
avec $P_b$ = 1013,25 hPa\\
et $\beta$ = -0,0065 K/m\\
et  $T_b$ = 288,15 K\\
et  $H_b$  = 0,0 m\\

Si H comprise entre 11000 et 20000 m \\
$P = P_b e^{\frac{-g_0}{R T}(H-H_b)}$ \\
avec $P_b$ = 226,320 hPa\\
et $\beta$ = 0 K/m\\
et  $T$ = 216,65 K\\
et  $H_b$  = 11000,0 m\\

Sinon si H sup�rieure �  20000 m  \\
 $P=misval$\\

